
\documentclass[x11names,svgnames]{beamer}

\usetheme{metropolis}
\setbeamertemplate{caption}{\raggedright\insertcaption\par}
\setbeamercolor{block title}{fg=LightSkyBlue4}
\newenvironment<>{importantblock}[1]{%
  \begin{actionenv}#2%
    \def\insertblocktitle{#1}%
    \par%
    \mode<presentation>{%
      \setbeamercolor{block title}{%
        use=normal text,
        fg=normal text.fg,
        bg=normal text.bg!80!fg
      }
      \setbeamercolor{block body}{
        use={block title, normal text},
        bg=block title.bg!50!normal text.bg
      }
    }%
    \usebeamertemplate{block begin}}
  {\par\usebeamertemplate{block end}\end{actionenv}}
\newenvironment<>{decisionblock}[1][Statement of the problem]{%
  \begin{actionenv}#2%
    \def\insertblocktitle{#1}%
    \par%
    \mode<presentation>{%
      \setbeamercolor{block title}{%
        use=normal text,
        fg=normal text.fg,
        bg=normal text.bg!80!fg
      }
      \setbeamercolor{block body}{
        use={block title, normal text},
        bg=block title.bg!50!normal text.bg
      }
    }%
    \usebeamertemplate{block begin}}
  {\par\usebeamertemplate{block end}\end{actionenv}}
\usepackage{fontspec}
\usepackage[french, frenchkw]{algorithm2e}
\SetKwFor{Pc}{Pour chaque}{faire}{fin}
\SetKwProg{Fn}{Fonction}{:}{}
\usepackage{textcomp}
\usepackage{minted}
\usepackage{pdfpages}
\usepackage{qrcode}
\usepackage[francais]{babel}
\usepackage{xspace}
\usepackage[normalem]{ulem}
\usepackage[backend=biber,style=authoryear,url=false]{biblatex}
\addbibresource{Bibliography.bib} % Bibliography file
\usepackage{amsmath}
\usepackage{mathtools}
\usepackage{pifont}% http://ctan.org/pkg/pifont
\newcommand{\cmark}{\ding{51}}%
\newcommand{\xmark}{\ding{55}}%
\usepackage{marvosym}
\usepackage{tikz}
\usetikzlibrary{calc, positioning, arrows, automata, fadings, decorations.pathmorphing, shapes, patterns, tikzmark, fit, chains, snakes, shapes.multipart}
% https://tex.stackexchange.com/questions/401884/how-do-i-change-hyperlinks-color-only
\hypersetup{
  colorlinks,
  allcolors=.,
  urlcolor=DarkBlue,
}


\author{}
\date{Lundi 7 octobre 2024}

\title{\includegraphics[width=0.75\textwidth]{../../img/logo-igm.png} \\ Algorithmique et Programmation 1}
\institute{L1 Mathématiques - L1 Informatique \\ Semestre 1}

\begin{document}

\maketitle

% Objectifs : Fonctions complexes: mémoire, prototypes/signature Appels : évaluation, retour, pile. Terminal. Exemple : recherche du minimum dans une liste

\begin{frame}{CC0 : Comment et quoi ?}
  \begin{block}{Horaire}
    \vspace{-0.25em}
    \begin{itemize}
    \item \textbf{Jeudi 10 octobre 10h45-11h15} (tiers-temps : 11h25) en A1
    \item Les retardataires ne seront pas admis
    \item Correction de 11h45 à 12h45 (en A1)
    \end{itemize}
  \end{block}
  \pause
  \begin{block}{Contenu}
    \vspace{-0.25em}
    \begin{itemize}
    \item Tout ce qu'on a vu, y compris la séance d'aujourd'hui
    \item Les slides \alert{et} les notes de cours
    \item Pas de pseudo-code
    \item QCM
    \end{itemize}
  \end{block}
  \pause
  \begin{block}{Sur la triche}
    \vspace{-0.25em}
    \begin{itemize}
    \item L'examen est \textbf{individuel}
    \item Tolérance zéro aux CC 1 et 2 (et 0.0000000000001 au CC0)
    \end{itemize}
  \end{block}
\end{frame}

\section{Approfondissements : les fonctions}

\begin{frame}{Fonction}
    En informatique, une fonction est :
    \begin{itemize}
    \item Un morceau de programme
    \item Portant en général \emph{un nom}
    \item Prenant un ou plusieurs \emph{paramètres} (ou zéro)
    \item Renvoyant un résultat (la plupart du temps)
    \end{itemize}
\end{frame}

\begin{frame}{Anatomie d'une fonction}
    \begin{tikzpicture}
      % Rectangle pour représenter la fonction
      \node [draw, rectangle, text width=2cm, minimum height=1.5cm, align=center] (function) at (0,0) {Fonction};

      % Chemin pour représenter l'entrée (paramètres)
      \draw[-latex] (-5,0) -- (-1.2,0) node[midway, above, label={below: (paramètres)}] {Entrée};

      % Chemin pour représenter la sortie (valeur de retour)
      \draw[-latex] (1.2,0) -- (5,0) node[midway, above, label= {below: (valeur de retour)}] {Sortie};

      % Flèche d'affichage (zig-zag)
      \onslide<2->{\draw [decorate, -latex, decoration={snake, segment length=5mm, amplitude=1mm, post length=1mm}] (0,0.5) -- (0,2.5) node[midway, left] {Affichages};}

      % Flèche d'effets de bord (zig-zag)
      \onslide<3->{\draw [decorate, -latex, decoration={snake, segment length=5mm, amplitude=1mm, post length=1mm}] (0,-0.5) -- (0,-2.5) node[midway, left, label={right: (aka: Effets secondaires)}] {Effets de bord};}
    \end{tikzpicture}
\end{frame}


\section{Exemples}

\begin{frame}[fragile]{Fonction à paramètres et résultats}
  Calculer le \emph{minimum} de deux entiers (dans Thonny) :
  \pause
\begin{minted}{python}
def minimum(a,b):
    if a < b:
        return a
    else:
        return b
\end{minted}
\end{frame}

\begin{frame}[fragile]{Fonction à paramètres et résultats}
  Calculer le \emph{minimum} d'une liste non-vide (dans Thonny) :
  \pause
\begin{minted}{python}
def minimum(lst):
    m = lst[0] # /!\ Erreur si lst est vide
    for elt in lst:
        if elt < m:
        m = elt
    return m
\end{minted}
\end{frame}

\begin{frame}[fragile]{Fonction sans valeur de retour}
  \begin{columns}
    \begin{column}{0.7\textwidth}
  Dessiner un carré fait du caractère \mintinline{python}{caractere} (dans Thonny) :
\end{column}
\begin{column}{0.1\textwidth}
\begin{verbatim}
*****
*****
*****
*****
*****
\end{verbatim}
\end{column}
\end{columns}
\pause
\vspace{-0.5em}
\begin{overprint}
  \onslide<2>
\begin{minted}{python}
def dessine_carre(n, caractere):
    i = 0
    while i < n:
        j = 0
        while j < n:
            print(caractere, end = '')
            j += 1
        print('\n', end = '')
        i += 1
    return
\end{minted}
  \onslide<3>
\begin{minted}{python}
def dessine_carre(n, caractere):
    for i in range(0, n):
        for j in range(0, n):
            print(caractere, end = '')
        print('\n', end = '')
    return
\end{minted}
\end{overprint}
\end{frame}

\begin{frame}
  \frametitle{Composition de fonctions}

  \begin{block}{En maths}
    $$\forall x \in X, g \circ f(x) = g(f(x))$$
  \end{block}
  
  \begin{block}{En informatique}
    \vspace{0pt}
    On peut appeler une fonction dans une fonction ! \newline (et ainsi de suite)
  \end{block}
\end{frame}

\begin{frame}[fragile]{Dessiner un carré : variante}
  \begin{columns}
    \begin{column}{0.1\textwidth}
\begin{verbatim}
*****
*****
*****
*****
*****
\end{verbatim}
    \end{column}
    \begin{column}{0.7\textwidth}
      Carré de côté \mintinline{python}{n} = \mintinline{python}{n} lignes de longueur \mintinline{python}{n}
    \end{column}
\end{columns}
\bigskip
On peut donc décomposer ainsi :
  \begin{itemize}
  \item
\begin{minted}{python}
def dessine_ligne(n,car):
    """Dessine *une ligne* de longueur n
       composée du caractère car."""
\end{minted}
  \item
\begin{minted}{python}
def dessine_carre(n,car):
    """Dessine *un carré* de longueur n
       composé du caractère car."""
\end{minted}
  \end{itemize}
  \MVRightarrow{} Ce que ça donne dans Thonny
\end{frame}
\begin{frame}[fragile]{Dessiner un carré : variante}
\begin{minted}{python}
def dessine_ligne(n, caractere):
 """dessine une ligne de longueur n
    composée du caractère car"""
    for j in range(0, n):
        print(caractere, end = '')
    print('\n', end = '')
    return
\end{minted}
\end{frame}
\begin{frame}[fragile]{Dessiner un carré : variante}
\begin{minted}{python}
def dessine_carre(n, caractere):
 """dessine un carré de longueur n
    composé du caractère car"""
    for i in range(0, n):
        dessine_ligne(n,caractere)
    return
\end{minted}
\end{frame}

\begin{frame}[fragile]{Erreur fréquente : confusion \begin{tabular}{l}paramètre / saisie \\ retour / affichage\end{tabular}}

    \vspace{-0.25em}
  \begin{columns}
    \begin{column}{0.7\textwidth}
  \begin{block}{Paramètre / saisie}
    \vspace{-0.25em}
\begin{minted}{python}
def minimum(a,b):
    a = int(input()) # NON !
    b = int(input()) # NON !
    if a <= b:
        return a
    else:
        return b
\end{minted}
  \end{block}
\pause
    \vspace{-0.25em}
  \begin{block}{Retour / affichage}
    \vspace{-0.25em}
\begin{minted}{python}
def maximum(a,b):
    if a >= b:
        print(a) # NON !
    else:
        print(b) # NON !
\end{minted}
  \end{block}
    \end{column}
    \pause
    \begin{column}{0.3\textwidth}
      Dans Thonny
    \end{column}
    \end{columns}
\end{frame}

\begin{frame}
  \frametitle{Fonctions et espaces de nom}
  \begin{itemize}
  \item Paramètres et variables d'une fonction : \alert{indépendantes} des autres variables du programme
  \item N'existent plus une fois l'exécution de la fonction terminée
  \item On les appelle des variables \emph{locales}
\end{itemize}

$\Rightarrow{}$ On peut renommer les variables d'une fonction.

$\rightarrow{}$ Démonstration sur Thonny (\texttt{minimum.py}).

\end{frame}

\begin{frame}
  \frametitle{(Contre-)exemple : intervertir des variables}
  Dans Thonny : \texttt{echange.py}
  \pause
  \begin{itemize}
  \item changer les valeurs de \mintinline{python}{a} et \mintinline{python}{b} dans la fonction \emph{n'a pas d'effet} sur \mintinline{python}{x} et \mintinline{python}{y} dans le programme principal !
  \item la variable \mintinline{python}{temp} n'existe plus après l'exécution de la fonction
  \end{itemize}
\end{frame}

\begin{frame}
  \frametitle{Sous le capot : sémantique d'un appel}

  \begin{block}{Espace de noms}
    \vspace{0pt}
    Ensemble de noms (variables, fonctions) défini à un certain point d'un programme
  \end{block}

  \MVRightarrow{} lors d'un appel de fonction, création d'un espace de nom \emph{local} :
  \begin{itemize}
  \item Paramètres $\rightarrow$ leur valeur lors de l'appel
  \item Variables locales
  \end{itemize}

  \begin{block}{Exemple}
    \vspace{0pt}
    \texttt{minimum\_liste.py} dans \href{https://pythontutor.com/}{Python Tutor} et Thonny.
  \end{block}

\end{frame}

\begin{frame}
  \frametitle{Espaces de noms lors de l'exécution}
  À un point du programme, l'ensemble des noms (variables, fonctions) connus est constitué de plusieurs espaces imbriqués (du plus ancien au plus récent) :
  \begin{itemize}
  \item espace de noms prédéfini (\mintinline{python}{print}, \mintinline{python}{input}, \dots)
  \item espace de noms global : définis dans le programme principal
  \item espaces de noms locaux imbriqués dans l'\emph{ordre chronologique} des appels
    \begin{itemize}
    \item paramètres de l'appel
    \item variables locales à la fonction
    \end{itemize}
  \end{itemize}
\end{frame}

\begin{frame}[fragile]
  \frametitle{Pile d'appels}
  L'empilement des espaces de noms obéit à une politique de \emph{pile}

  \begin{itemize}
  \item sommet : appel en cours
  \item en-dessous : appels précédents
  \item (presque) tout en bas : espace de nom global
  \end{itemize}

  \begin{columns}
    \begin{column}{0.5\textwidth}
\begin{minted}{python}
def f(x):
    return g(x) + 1
def g(x):
    return h(x) + 2
def h(x):
    return x + 3
\end{minted}
    \end{column}
    \begin{column}{0.4\textwidth}
      \alt<1-4>{%
      \begin{tikzpicture}
        % Stack rectangle
        \draw (0,0) rectangle (2,3);
        % Stack top
        \node at (1,3.3) {Sommet};
        % Stack elements
        \pause
        \draw (0.2,0.5) rectangle (1.8,0.8) node[midway] {f};
        \pause
        \draw (0.2,0.9) rectangle (1.8,1.2) node[midway] {g};
        \pause
        \draw (0.2,1.3) rectangle (1.8,1.6) node[midway] {h};
        % Stack direction
        \draw[->,ultra thick] (2.5,0) -- (2.5,3);
      \end{tikzpicture}%
    }{%
      \begin{tikzpicture}[yscale=-1]
        % Stack rectangle
        \draw (0,0) rectangle (2,3);
        % Stack top
        \node at (1,3.3) {Sommet};
        % Stack elements
        \pause
        \draw (0.2,0.5) rectangle (1.8,0.8) node[midway] {f};
        \pause
        \draw (0.2,0.9) rectangle (1.8,1.2) node[midway] {g};
        \pause
        \draw (0.2,1.3) rectangle (1.8,1.6) node[midway] {h};
        % Stack direction
        \draw[->,ultra thick] (2.5,0) -- (2.5,3);
      \end{tikzpicture}%
    }
    \end{column}
  \end{columns}

  \medskip

  \onslide<5>{\textbf{Attention !} Dans Python Tutor, le plus récent est en bas.}
\end{frame}

\begin{frame}
  \frametitle{Pile d'appels}

  \begin{block}{Quand un nouvel appel commence}
    \begin{itemize}
    \item l'exécution de la fonction en cours s'interrompt
    \item un nouvel espace de noms local est créé
    \item l'exécution de la fonction appelée commence
    \end{itemize}
  \end{block}


  \begin{block}{Quand l'appel en cours se termine}
    \begin{itemize}
    \item son espace de nom est supprimé de la pile
    \item l'exécution de l'appel précédent reprend
    \end{itemize}

  \end{block}

\end{frame}
\begin{frame}{Portée des variables}

  \begin{block}{Accès à la valeur d'une variable}
  \begin{itemize}
  \item possible pour n'importe quel nom défini dans un des espaces de noms antérieurs
  \item si plusieurs espaces contiennent le même nom, c'est le plus récent qui est sélectionné
  \end{itemize}
\end{block}

\begin{block}{Affectation}
  
  \begin{itemize}
  \item par défaut, uniquement aux variables locales
  \item pour une variable dans un espace de nom plus ancien, mots-clés \mintinline{python}{global} ou \mintinline{python}{nonlocal} \newline (à utiliser avec précaution)
  \end{itemize}

\end{block}

\end{frame}

\begin{frame}[fragile]{Signature d'une fonction}

Considérons la fonction :
\begin{minted}{python}
def f(p_1, ..., p_n):
...
return expr
\end{minted}

Les noms \mintinline{python}{p_1} à \mintinline{python}{p_n} sont appelés paramètres (formels).

En général, ils d'un certain type. On peut le préciser en explicitant la \emph{signature} de la fonction :
\begin{minted}{python}
def f(p_1: type_1, ..., p_n: type_n) -> type_retour:
...
return expr
\end{minted}
\begin{exampleblock}{Exemple}
\begin{minted}{python}
def minimum(a: int, b: int) -> int:
\end{minted}
\end{exampleblock}
\end{frame}

\begin{frame}[fragile]{Déroulement détaillé d'un appel}
      Considérons maintenant l'appel \mintinline{python}{f(e_1, ..., e_n)} :
      \begin{itemize}
      \item les valeurs \mintinline{python}{v_1} à \mintinline{python}{v_n} sont appelées paramètres effectifs \newline (ou arguments)
      \item création d'un espace de noms local contenant \mintinline{python}{p_1} à \mintinline{python}{p_n} au sommet de la pile d'appels
      \item chaque expression \mintinline{python}{e_i} est évaluée en une valeur \mintinline{python}{v_i} et affectée à la variable \mintinline{python}{p_i}
      \item exécution du corps de la fonction dans l'espace de noms local
        \pause
      \item si la fonction exécute l'instruction \mintinline{python}{return expr} ou atteint la fin de son bloc d'instructions :
        \begin{itemize}
        \item l'espace de noms local est détruit
        \item l'expression appelante \mintinline{python}{f(e_1, ..., e_n)} prend la valeur de \mintinline{python}{expr} (respectivement \mintinline{python}{None})
        \item reprise du programme principal dans l'espace global
        \end{itemize}
      \end{itemize}

\end{frame}
\begin{frame}[fragile]{Documentation et test de fonctions}
  \begin{block}{Chaînes de documentation (\mintinline{python}{docstring})}
    \vspace{0pt}

    Bonne pratique : indiquer par un commentaire
    \vspace{-0.5em}
    \begin{itemize}
    \item à quoi sert une fonction
     \item ce que représentent ses paramètres et leur type
     \item ce que représente sa valeur de retour
     \item d'éventuels effets ou causes secondaires
    \end{itemize}
  \end{block}

\begin{minted}{python}
def triple(n):
  """
  Calcule le triple du nombre n (int ou float)
  ou la répétition trois fois de la chaîne n.
  """
  return n * 3
\end{minted}
\end{frame}
\begin{frame}[fragile]{Documentation et test de fonctions}
  \begin{itemize}
    \item On peut accéder à la chaîne de documentation d'une fonction en tapant \mintinline{python}{help(nom de la fonction)} dans l'interpréteur
  \item Cela fonctionne aussi pour les fonctions prédéfinies ou issues de modules
  \end{itemize}
\end{frame}
\begin{frame}[fragile]{Tests intégrés à la documentation (\mintinline{python}{doctest})}
  \begin{itemize}
  \item Toute fonction \emph{doit être testée immédiatement} pour s'assurer qu'elle fonctionne.
  \item On peut intégrer les tests à sa documentation.
  \end{itemize}
\begin{minted}{python}
def triple(n):
  """
  Calcule le triple du nombre n (int ou float)
  ou la répétition trois fois de la chaîne n.
  
  >>> triple(3)
  9
  >>> triple(9.0)
  27.0
  """
  return n * 3
\end{minted}
\end{frame}
\begin{frame}[fragile,fragile]{Tests intégrés à la documentation (\mintinline{python}{doctest})}
  Il existe des outils qui permettent de lancer automatiquement tous les tests présents dans la documentation, et de vérifier qu'ils produisent les résultats annoncés.

  Par exemple, à la fin d'un programme, on peut écrire le code suivant pour lancer systématiquement tous les tests présents dans le fichier :
\begin{minted}{python}
  import doctest
  doctest.testmod()
\end{minted}
\end{frame}

\end{document}

%%% Local Variables:
%%% TeX-engine: xetex
%%% TeX-command-extra-options: "-shell-escape"
%%% End:
